\documentclass[11pt]{article}
\usepackage[T1]{fontenc}
\usepackage[utf8]{inputenc}
\usepackage[english]{babel}
\usepackage[margin=1in]{geometry}
\usepackage{amsmath,amssymb}
\usepackage{enumitem}
\usepackage{microtype}
\usepackage[hidelinks]{hyperref}

\setlength{\parindent}{0pt}
\setlength{\parskip}{0.65em}
\setlist{leftmargin=*,topsep=0.3em,itemsep=0.35em}

\title{Technical Review of\\
\emph{Early-Stopping Activation Steering for Factual Preference Alignment in Vietnamese Medical Language Models}}
\author{}
\date{September 3, 2026}

\begin{document}
\maketitle

\section*{Overall assessment}

The revision corrects the Qwen technical-report identifier, repairs the author metadata for Refs.~\texttt{ref2} and \texttt{ref9}, and removes the previous overfull Hybrid RAG table.  It also places the correct Continuous-versus-baseline difference, $\Delta=+2.95$, in contribution 3.  The exact McNemar $p$-values and the intervention-dose algebra remain internally coherent.

Major revision is nevertheless required.  The corrected norm difference appears only in the contribution list; the Results section still reports the old incorrect values and conflicts with the claimed fixed teacher-forcing design.  The directional-control paragraph has become syntactically corrupted around a missing item~(iii), and its pair-shuffle and partial label-shuffle controls remain scientifically inadequate.  The token-timing convention, T4/\texttt{bfloat16} environment, benchmark validity, RefPref interpretation, generation-cap effects, statistical interval labels, and unsupported Experiment~10 claims are still unresolved.  The strongest completed-answer result remains Hard Cutoff, whereas the paper's main narrative emphasizes Linear Decay.

\section*{Major technical findings}

\subsection*{1. The norm correction is incomplete and introduces a new unexplained quantity}

\textbf{Classification: definite internal inconsistency and arithmetic error.}

\textbf{Location.} Abstract; contribution 3; Empirical Activation-Norm Dynamics; Discussion.

\textbf{Problem.} Contribution 3 now correctly states that the Continuous value at $t=50$ is $56.34-53.39=+2.95$ above baseline.  The Results paragraph still calls the elevation ``$+1.97$'' and still describes Linear Decay at $t=50$ as a ``norm undershoot of $-3.01$ units relative to baseline $53.39$,'' even though $53.39-53.39=0$.  The complete displayed differences are
\begin{align*}
56.28-53.64 &= +2.64 \quad (t=10,\ \text{Continuous}),\\
56.34-53.39 &= +2.95 \quad (t=50,\ \text{Continuous}),\\
52.95-53.64 &= -0.69 \quad (t=10,\ \text{Linear Decay}),\\
53.39-53.39 &= 0 \quad (t=50,\ \text{Linear Decay}),\\
52.75-53.39 &= -0.64 \quad (t=50,\ \text{Hard Cutoff}).
\end{align*}
Contribution 3 additionally introduces ``$53.39\pm0.64$'' without defining whether $0.64$ is a standard deviation, confidence interval, range, or the absolute Hard Cutoff difference.  No uncertainty statistic of 0.64 is reported in Results.

The design contradiction also remains.  Contribution 3 calls the analysis fixed teacher forcing.  If identical tokens are teacher-forced and $\alpha(t)=0$ for $t>K$, Linear Decay and Hard Cutoff must both equal baseline at the measured layer at $t=50$.  Linear Decay does, but Hard Cutoff does not.  If the values instead use free-running trajectories, their post-$K$ differences reflect different token histories and cannot isolate direct hook relaxation.

\textbf{Why it matters.} Activation-norm mitigation is the proposed mechanism and a named contribution.  Contradictory values, undefined uncertainty, and incompatible evaluation designs prevent verification.

\textbf{Suggested fix.} Update every manuscript occurrence from one documented run.  Define all uncertainty terms and report pre-hook and post-hook norms, $\|\hat h-h\|_2$, $\langle h,v_{\text{steer}}\rangle$, sample counts, uncertainty, and complete trajectories.  State unambiguously whether tokens are fixed or free-running.  Under teacher forcing, present equality after $K$ as a direct consequence of $\alpha(t)=0$, not separate behavioral evidence.

\subsection*{2. The directional-control specification is malformed, and key null controls remain invalid}

\textbf{Classification: definite source defect, definite control-design error, and implementation ambiguity.}

\textbf{Location.} Contrastive Vector Estimation \& Directional Controls; Table~\texttt{tab:baselines}; directional-specificity claims.

\textbf{Problem.} The revised five-condition sentence proceeds from item~(ii) directly to an unlabeled phrase after the text
\begin{quote}
``$\cos(v_{\text{steer}},v_{\text{shuf}})=0.5839$)), permuting pair indices $\pi(i)$''
\end{quote}
and then resumes at item~(iv).  The label ``(iii) Pair-Shuffled Steering ($v_{\text{pair\_shuf}}$)'' has been deleted, and an extra closing parenthesis remains.  The method therefore no longer parses as a complete control list.

Even when restored, pair shuffling is invariant for the stated estimator:
\[
\frac{1}{N}\sum_i\left(h_i^+-h_{\pi(i)}^-\right)=\mu^+-\mu^-.
\]
Its cosine of 1.0000 is mathematically inevitable, so it is not a negative control.  Label shuffling remains unclear because $N_{\text{flipped}}=185$ is not reconciled with the stated 10,290 training pairs.  Flipping only 185 of 10,290 examples leaves almost all labels unchanged; if a smaller vector-construction subset was used, it is not documented.

The covariance control also uses undefined $z_r$ in $v_{\text{cov}}^{(r)}=Lz_r/\|Lz_r\|_2$.  The normalized isotropic direction is written as Gaussian even though $z_r\sim\mathcal N(0,I)$ followed by $v_r=z_r/\|z_r\|_2$ is a unit-sphere distribution.  The coefficient and temporal schedule used for the shuffled and random controls are not explicitly stated.

Finally, isotropic directions average 55.82\% RefPref versus a 73.00\% unsteered baseline, showing that large arbitrary perturbations are broadly damaging.  Beating this harmful distribution does not by itself establish truthfulness specificity.  The learned vector is only 2.39~pp above the covariance-matched mean, with no vector-level range, exceedance count, or paired null test.

\textbf{Why it matters.} Directional specificity is a central contribution, but the current controls conflate truth-related structure, response-style artifacts, and generic perturbation damage.

\textbf{Suggested fix.} Repair the control list and remove Pair-Shuffled Steering as evidence.  Use full balanced label permutations or random signs on pairwise differences over the complete documented training set.  Define every random variable and match layer, norm, coefficient, schedule, and prompts across controls.  Report each vector's paired change from baseline and a pre-specified empirical tail test, and add length-, EOS-, and final-token-matched extraction controls.

\subsection*{3. Token timing, layer terminology, and the negative-steering convention remain implementation-critical}

\textbf{Classification: likely off-by-one implementation error and notation/sign ambiguity.}

\textbf{Location.} Early-Stopping Steering Hook; Experimental Setup; abstract; Directional Controls and Results.

\textbf{Problem.} Prompt prefill is explicitly unsteered.  In standard cached autoregressive generation, prefill logits produce the first output token; the first post-prefill call consumes that token and produces logits for the second.  The described $t=1$ intervention therefore affects predicted output tokens 2 through $K+1$, not the ``first $K$ generated tokens,'' unless a custom loop applies the hook differently.

The manuscript alternates between ``Layer~8,'' $h_8$, and \texttt{model.model.layers[8]}, which is correctly identified as the ninth decoder block.  The scientific layer label and code index should not be used interchangeably.

The revised negative-control equation, $\hat h_8^{(t)}=h_8^{(t)}-|\alpha(t)|v_{\text{steer}}$, is negative steering.  However, the same condition is still described as both ``$-v_{\text{steer}}$'' and ``$\alpha_0=-18.0$.''  An implementation based on the generic equation $h+\alpha v$ must use one sign convention, while an implementation using the explicit subtraction must not negate the vector again.  The source does not show which code path generated the results.

The extraction state also depends on the final token of $x_i\oplus y_i$, but the chat template, EOS inclusion, response-length handling, answer separator, and equivalence of extraction and injection sites are not specified.  The dose formulas additionally assume $T\ge K$; shorter sequences require an upper limit of $\min(T,K)$.

\textbf{Why it matters.} These choices determine which output logits are affected, which decoder block is modified, and the direction of a key experimental condition.

\textbf{Suggested fix.} Provide executable pseudocode for prefill and cached decoding, naming both the consumed token and the predicted token.  Use ``decoder block at zero-based index 8 (ninth block)'' consistently.  Implement negative steering either as $(v,\alpha)=(v_{\text{steer}},-18)$ in the generic equation or as explicit subtraction with a positive magnitude, not a mixture of both descriptions.  Specify tokenization, chat template, EOS handling, tensor tuple reconstruction, and the $T<K$ cases.

\subsection*{4. The Tesla T4 and \texttt{bfloat16} environment is not self-consistent}

\textbf{Classification: externally verified likely implementation/reporting error.}

\textbf{Location.} Model Checkpoint \& Precision; Reproducibility Environment; latency claims.

\textbf{Problem.} The manuscript reports NF4 with \texttt{bfloat16} compute on NVIDIA Tesla T4 GPUs.  Tesla T4 is a Turing, compute-capability-7.5 device, whereas native CUDA \texttt{bfloat16} requires compute capability 8.0 or higher.  NVIDIA's T4 specifications list FP16/FP32 mixed precision, INT8, and INT4, but not native BF16.  The runs therefore likely used FP16, a fallback/conversion path, different hardware, or incorrectly reported metadata.

\textbf{Why it matters.} Actual compute dtype affects numerical outputs, memory, kernel selection, and the latency comparisons used as a systems contribution.

\textbf{Suggested fix.} Report the outputs of the PyTorch CUDA checks \texttt{get\_device\_capability} and \texttt{is\_bf16\_supported}, together with actual parameter, activation, autocast, and BitsAndBytes dtypes from the run logs.  If T4 was used, report FP16 unless a documented fallback was used; if BF16 ran natively, correct the GPU model.  Re-run timing after reconciling the environment.

\subsection*{5. Benchmark documentation and RefPref do not support clinical factual-accuracy language}

\textbf{Classification: unsupported claim and likely dataset-validity issue.}

\textbf{Location.} Benchmark Construction; BERTScore Reference-Preference Criterion; abstract, Results, table captions, and conclusion.

\textbf{Problem.} The manuscript does not document expert qualifications, source-to-question conversion, synthetic-negative prompts/model, factual review, agreement, adjudication, or correction procedures.  The category counts 165, 160, and 175 sum to 500 and therefore describe the evaluation subset, but they are introduced as though they characterize all 14,700 records.  The selection of 500 questions from the 2,205-example test partition is unexplained.  ``Question-disjoint'' is not defined at the drug, entity, or formulary-passage level, so source-level overlap may remain between vector estimation and testing.

RefPref only identifies whether an output is semantically closer to one reference than to one synthetic counter-answer.  It does not verify dosage arithmetic, contraindications, interactions, or additional unsupported medical statements.  It is also sensitive to answer length and style.  Although the limitations section calls it a proxy, the abstract and conclusion still switch its descriptor to ``factual accuracy,'' ``factual alignment,'' or clinical-safety language.  Tie counts and complete BERTScore options, including IDF and baseline rescaling, are omitted.

\textbf{Why it matters.} A semantic preference proxy is not a clinical correctness or safety endpoint.  Synthetic-answer artifacts or source leakage could contribute materially to the measured effect.

\textbf{Suggested fix.} Document the full data-generation and adjudication pipeline, split-level category counts, sampling rule, and source/drug overlap checks.  Group splits by source passage and drug/entity and release immutable split IDs.  Add blinded clinician adjudication and structured scoring for dose, interaction, contraindication, and unsupported-claim errors.  Report ties and all BERTScore settings, and call the current endpoint ``RefPref rate'' consistently.

\subsection*{6. The schedule narrative is confounded with generation truncation}

\textbf{Classification: claim/evidence mismatch and terminology drift.}

\textbf{Location.} Abstract; natural-completion and bounded evaluations; factorial ablation; conclusion.

\textbf{Problem.} Under natural completion, Hard Cutoff is the only schedule with a Holm-adjusted significant improvement.  Linear Decay gains 2.40~pp with $p_{\text{adj}}=0.1189$.  The bounded experiment instead promotes Linear Decay, but the ablation reports 0.0\% EOS for that condition and at most 0.2\% for any condition.  Thus nearly every 200-token response is a truncated prefix.  The unsteered baseline changes from 73.00\% at 200 tokens to 65.40\% under natural completion, directly showing cap sensitivity.

The factorial table does not report RefPref, the paper's principal endpoint, and provides no paired uncertainty for schedule comparisons.  It therefore cannot establish which schedule is best for reference preference.  ``Early-stopping steering'' is used sometimes for Linear Decay alone and elsewhere as an umbrella including Hard Cutoff.

\textbf{Why it matters.} Evidence on truncated prefixes cannot be assumed to transfer to completed answers, and the completed-answer test currently supports a different schedule from the main Linear Decay narrative.

\textbf{Suggested fix.} Pre-specify a primary schedule and endpoint.  Rename the 200-token experiment as truncated-prefix stress testing, report output-length distributions, and give paired RefPref comparisons for all schedules at both caps.  State that current natural-completion evidence favors Hard Cutoff and that Linear Decay is exploratory there.  Use schedule-specific names consistently.

\subsection*{7. The degeneration mechanism and Experiment 10 remain unsupported}

\textbf{Classification: unsupported causal claim and unsupported result.}

\textbf{Location.} Introduction; natural-completion Results; Activation-Norm Dynamics; Multi-Metric Pareto Trade-off \& Deployment Mitigation.

\textbf{Problem.} Continuous steering is said to produce repetitive loops and degraded syntax, but no syntax metric or loop analysis is provided.  The only completed-answer repetition metric is lower for Continuous Steering (4.38\%) than for Hard Cutoff (5.35\%) or Linear Decay (5.37\%).  EOS occurrence also does not prove the absence of loops.

The Discussion cites ``Experiment 10'' and reports $\theta_{\text{rep}}=1.15$, Rep-4 of 3.76\%, and latencies of 78.74~s and 45.05~s.  No Experiment~10 method, table, test population, output length, uncertainty, or RefPref result appears in the paper.  Consequently, the statements that the trade-off is ``resolved completely'' and that the $+5.20$~pp gain is preserved are unsupported.  The same paragraph later calls the penalty only a ``viable path,'' contradicting the claimed completed experiment.

\textbf{Why it matters.} The paper's mechanistic motivation and deployment remedy are not supported by the presented behavioral evidence.

\textbf{Suggested fix.} Remove these claims or add a complete matched experiment covering every schedule with RefPref, BERTScore, Rep-4, syntax/degeneration measures, EOS, generated-token counts, and synchronized latency.  Until then, describe the relationship between norm offset and degeneration as an unresolved hypothesis.

\subsection*{8. Confidence-interval labels and endpoint descriptions remain inconsistent}

\textbf{Classification: likely statistical-reporting error; exact tests are internally correct.}

\textbf{Location.} Contribution 5; Tables~\texttt{tab:mcnemar} and \texttt{tab:mcnemar\_bounded}; Hybrid and BM25 comparisons.

\textbf{Problem.} The displayed exact two-sided McNemar $p$-values reproduce from the discordant counts.  However, $[+1.37,+7.03]$~pp and $[+2.25,+8.15]$~pp numerically equal the paired normal intervals
\[
\widehat\Delta\pm1.96\sqrt{\frac{b+c-(c-b)^2/N}{N^2}},
\]
while Table~\texttt{tab:mcnemar} labels its intervals as a $B=10{,}000$ percentile bootstrap with seed 42.  Direct paired percentile resampling of the displayed four-cell counts does not reproduce those endpoints under common seeded implementations.  No per-example data or code is provided to resolve the discrepancy.

The Wilcoxon test concerns continuous BERTScore F1, not binary RefPref, yet contribution 5 says it confirms reference-preference gains.  In the Hybrid RAG table, the prose calls the discordant cells $b=22,c=27$, whereas the displayed row-by-column layout would conventionally label 27 as the upper-right $b$ cell and 22 as $c$.  The two-sided test is unaffected, but the notation changes orientation across tables.

\textbf{Why it matters.} Statistical validation is a named contribution.  Mislabeling interval methods and endpoints prevents reliable interpretation and reproduction.

\textbf{Suggested fix.} Release paired binary outcomes, per-example BERTScore values, and exact scripts.  Name every interval method accurately, define the comparison family for Holm adjustment, specify Wilcoxon zero handling, and keep one $b,c$ orientation throughout.

\subsection*{9. The Hybrid RAG result does not support an ``outperforms'' claim}

\textbf{Classification: unsupported comparative claim and reproducibility ambiguity.}

\textbf{Location.} RAG Results; Tables~\texttt{tab:baselines} and \texttt{tab:mcnemar\_rag}; abstract and conclusion.

\textbf{Problem.} Linear Decay is said to ``outperform Hybrid RAG'' by 1.00~pp, but the same section reports exact paired $p=0.5682$ and correctly calls the RefPref rates statistically indistinguishable.  Hybrid RAG also has the higher reference BERTScore F1, 0.7190 versus 0.7150.  Steering offers lower reported context and latency, but it does not dominate Hybrid RAG across the reported metrics.

The retrieval setup omits passage construction, retrieval top-$k$, dense pooling and normalization, model revision, RRF inputs, context selection/truncation, prompt templates, and the \texttt{pyvi} version.  It is unclear whether reported latency includes retrieval, CUDA synchronization and warm-up, and whether output lengths were matched.  Identical $7.32\pm0.84$~s entries for every non-RAG condition do not independently establish zero hook overhead.

\textbf{Why it matters.} The best non-oracle RAG comparator is statistically tied on RefPref and better on absolute BERTScore.  The useful systems trade-off should not be converted into generic performance superiority.

\textbf{Suggested fix.} Replace ``outperforms'' with ``numerically exceeds'' and present a quality--latency--context trade-off.  Fully specify retrieval and timing, report output-token distributions, decompose retrieval/prefill/decoding time, and provide paired uncertainty for latency.  Keep the BM25 claim limited to the specific weak lexical baseline.

\subsection*{10. The layer-sweep interpretation exceeds the reported evidence}

\textbf{Classification: unsupported inference.}

\textbf{Location.} Layer-Wise Subspace Probing.

\textbf{Problem.} A BERTScore range of 0.7040--0.7140 over layers 4--15 on $N_{\text{eval}}=100$, without individual results or paired uncertainty, does not demonstrate that truthfulness representations are distributed across layers and does not ``confirm framework robustness.''  Similar scores may arise from residual propagation, layer correlation, or sampling noise.

\textbf{Why it matters.} An exploratory selection sweep is being presented as a representation-level and robustness conclusion.

\textbf{Suggested fix.} Report all layer results, paired intervals, the selection criterion, and the complete hyperparameter search.  Limit the conclusion to similar validation scores among several tested intervention layers unless a dedicated localization analysis is added.

\section*{Reference integrity and meaningful editorial findings}

The revision correctly updates Ref.~\texttt{ref4} to arXiv:2412.15115 and repairs Ref.~\texttt{ref9}.  The following verified defects remain:

\begin{enumerate}
 \item \textbf{Ref.~\texttt{ref1}:} the seven-author list still does not match the official author list for arXiv:2310.01405.  Re-export the record from the primary source.
 \item \textbf{Ref.~\texttt{ref2}:} the author list is now correct, but the NeurIPS page range is not.  Replace 42910--42938 with 41451--41530.
 \item \textbf{Ref.~\texttt{ref8}:} the current arXiv:2308.10248 record is titled \emph{Steering Language Models With Activation Engineering} and lists Alexander Matt Turner, Lisa Thiergart, Gavin Leech, David Udell, Juan J.~Vazquez, Ulisse Mini, and Monte MacDiarmid.  The manuscript uses an older title, changes the author order, and omits Vazquez; re-export the intended version consistently.
\end{enumerate}

The Reproducibility Environment sentence remains duplicated in consecutive paragraphs.  Keep one copy and merge the seed information into it.  The Discussion also repeats the natural-completion trade-off within the same paragraph.

\section*{Proofing sweep}

The bundled scan was run on the revised TeX source and freshly compiled PDF.  Its double-period hit is the valid ellipsis in $\{1,\ldots,100\}$, and its semicolon-capitalization hit is the proper metric name ``Recall@3''; neither is an error.  No inverse-trigonometric or coordinate-transform expression occurs, so no branch or quadrant ambiguity was found.

High-confidence corrections are:

\begin{itemize}
 \item \textbf{Directional Controls:} restore the missing ``(iii) Pair-Shuffled Steering ($v_{\text{pair\_shuf}}$)'' label if the condition is retained, and delete the duplicated closing parenthesis after 0.5839.  Scientifically, the invariant control should then be removed or replaced as described above.
 \item \textbf{Natural-completion Results:} 5.35\% versus 4.10\% is a 30.5\% increase, while 5.37\% versus 4.10\% is approximately 31.0\%.  Do not apply 30.5\% to both schedules.
 \item \textbf{RAG Results:} the Hybrid RRF sentence repeats ``achieving'' in the same clause.  Replace the second occurrence with ``and yields.''
 \item \textbf{Units:} write $7.32$~s and $+7.13$~s consistently rather than \texttt{7.32s} and \texttt{+7.13s}.
 \item \textbf{Control organization:} ``we evaluate two controls'' is followed by a second item claiming five conditions and repeating Negative Steering.  Replace the nested structure with one flat list matching the results table.
\end{itemize}

\section*{Highest-priority revision sequence}

\begin{enumerate}
 \item Reconcile all norm values, uncertainty, and teacher-forcing/free-generation descriptions from one experiment.
 \item Repair and redesign the directional controls; specify token timing, layer indexing, sign convention, and vector extraction in executable pseudocode.
 \item Reconcile T4 hardware with the actual compute dtype and repeat timing measurements if necessary.
 \item Document benchmark construction and source-level splitting, then add clinician-adjudicated factual evaluation.
 \item Separate completed-answer evidence from truncated-prefix evidence and pre-specify the primary schedule.
 \item Remove or fully report Experiment~10, correct interval labels, and recast the Hybrid RAG comparison.
 \item Finish bibliography and proofing corrections.
\end{enumerate}

\section*{Items requiring external verification or unavailable artifacts}

The following cannot be established from the manuscript alone:

\begin{itemize}
 \item expert qualifications, factual adjudication, availability, and redistribution permission for the benchmark and formulary-derived passages;
 \item whether Ref.~\texttt{ref3} supports the specific claims about dosage, interaction, and contraindication errors, rather than only broader clinical limitations;
 \item whether Ref.~\texttt{ref11} supports the asserted internal-activation explanation for ignoring retrieved evidence;
 \item run logs showing GPU identity, CUDA capability, actual compute dtype, fallback behavior, quantization configuration, and model revision;
 \item source code, immutable split IDs, retrieval index, per-query outputs, and statistical scripts needed to reproduce norm, latency, control, RAG, Wilcoxon, and interval claims.
\end{itemize}

\section*{Checks that were internally consistent}

The 70/15/15 split counts sum to 14,700; the three evaluation categories sum to 500; category counts reproduce the stated overall RefPref totals; $D_{\text{decay}}=(K+1)|\alpha_0|/2$ and $\alpha_{\text{match}}=0.0425\alpha_0$ are correct for $K=16$, $T=200$, and $T\ge K$; and the displayed exact McNemar $p$-values agree with the discordant counts.  The official Qwen2.5-7B-Instruct configuration supports 28 decoder blocks and hidden size 3584.  No dimensional error was identified in the schedule equations, and no inverse-trigonometric branch issue is present.  These clean checks do not resolve the norm, timing, hardware, control, and endpoint problems above.

\end{document}